% -----------------------------------------------------------
\section{Your life}
% -----------------------------------------------------------
	\label{GM-LIFE}
	
	% ----------------------
	\subsection{Living the life of your dreams} % *DREAMS
	% ----------------------
		\label{GM-LIFE-DREAMS}
		
		% -----
		\subsubsection{0001 | Peter Keating on the courage to live the life you want}
		% -----
			\label{GM-LIFE-DREAMS-0001}
		
			Katie...for six years...I thought of how I'd ask your forgiveness some day. And now I have the chance, but I won't ask it. It seems...it seems beside the point. I know it's horrible to say that, but that's how it seems to me. It was the worst thing I ever did in my life---but not because I hurt you. I did hurt you, Katie, and maybe more than you know yourself. But that's not my worst guilt. ...Katie, I wanted to marry you. It was the only thing I ever really wanted. And that's the sin that can't be forgiven---that I hadn't done what I wanted. It feels so dirty and pointless and monstrous, as one feels about insanity, because there's no sense to it, no dignity, nothing but pain---and wasted pain. ...Katie, why do they always teach us that it's easy and evil to do what we want and that we need discipline to restrain ourselves? It's the hardest thing in the world---to do what we want. And it takes the greatest kind of courage. I mean, what we really want. As I wanted to marry you. Not as I want to sleep with some woman or get drunk or get my name in the papers. Those things---they're not even desires---they're things people do to escape from desires---because it's such a big responsibility, really to want something. (\textit{The Fountainhead}, 598--599)
			
		% -----
		\subsubsection{0002 | John Galt on reaching the life you envision}
		% -----
			\label{GM-LIFE-DREAMS-0002}
		
			In that world, you'll be able to rise in the morning with the spirit you had known in your childhood: that spirit of eagerness, adventure and certainty which comes from dealing with a rational universe. No child is afraid of nature; it is your fear of men that will vanish, the fear that has stunted your soul, the fear you acquired in your early encounters with the incomprehensible, the unpredictable, the contradictory, the arbitrary, the hidden, the faked, the irrational in men. You will live in a world of responsible beings, who will be as consistent and reliable as facts; the guarantee of their character will be a system of existence where objective reality is the standard and the judge. Your virtues will be given protection, your vices and weaknesses will not. Every chance will be open to your good, none will be provided for your evil. What you'll receive from men will not be alms, or pity, or mercy, or forgiveness of sins, but a single value: justice. And when you'll look at men or at yourself, you will feel, not disgust, suspicion and guilt, but a single constant: respect.

			Such is the future you are capable of winning. It requires a struggle; so does any human value. All life is a purposeful struggle, and your only choice is the choice of a goal. Do you wish to continue the battle of your present or do you wish to fight for my world? Do you wish to continue a struggle that consists of clinging to precarious ledges in a sliding descent to the abyss, a struggle where the hardships you endure are irreversible and the victories you win bring you closer to destruction? Or do you wish to undertake a struggle that consists of rising from ledge to ledge in a steady ascent to the top, a struggle where the hardships are investments in your future, and the victories bring you irreversibly closer to the world of your moral ideal, and should you die without reaching full sunlight, you will die on a level touched by its rays? Such is the choice before you. Let your mind and your love of existence decide.

			The last of my words will be addressed to those heroes who might still be hidden in the world, those who are held prisoner, not by their evasions, but by their virtues and their desperate courage. My brothers in spirit, check on your virtues and on the nature of the enemies you're serving. Your destroyers hold you by means of your endurance, your generosity, your innocence, your love---the endurance that carries their burdens---the generosity that responds to their cries of despair---the innocence that is unable to conceive of their evil and gives them the benefit of every doubt, refusing to condemn them without understanding and incapable of understanding such motives as theirs---the love, your love of life, which makes you believe that they are men and that they love it, too. But the world of today is the world they wanted; life is the object of their hatred. Leave them to the death they worship. In the name of your magnificent devotion to this earth, leave them, don't exhaust the greatness of your soul on achieving the triumph of the evil of theirs. Do you hear me...my love?

			In the name of the best within you, do not sacrifice this world to those who are its worst. In the name of the values that keep you alive, do not let your vision of man be distorted by the ugly, the cowardly, the mindless in those who have never achieved his title. Do not lose your knowledge that man's proper estate is an upright posture, an intransigent mind and a step that travels unlimited roads. Do not let your fire go out, spark by irreplaceable spark, in the hopeless swamps of the approximate, the not-quite, the not-yet, the not-at-all.

			Do not let the hero in your soul perish, in lonely frustration for the life you deserved, but have never been able to reach. Check your road and the nature of your battle. The world you desired can be won, it exists, it is real, it is possible, it's yours.

			But to win it requires your total dedication and a total break with the world of your past, with the doctrine that man is a sacrificial animal who exists for the pleasure of others. Fight for the value of your person. Fight for the virtue of your pride. Fight for the essence of that which is man: for his sovereign rational mind. Fight with the radiant certainty and the absolute rectitude of knowing that yours is the Morality of Life and that yours is the battle for any achievement, any value, any grandeur, any goodness, any joy that has ever existed on this earth.

			You will win when you are ready to pronounce the oath I have taken at the start of my battle---and for those who wish to know the day of my return, I shall now repeat it to the hearing of the world: ``I swear---by my life and my love of it---that I will never live for the sake of another man, nor ask another man to live for mine.'' (\textit{Atlas Shrugged})
				
	% ----------------------
	\subsection{Energy, youth, and power} % *ENERGY
	% ----------------------
		\label{GM-LIFE-ENERGY}
		
		% -----
		\subsubsection{0001 | Pook's ``Fountain of Youth''} % *FOUNTAINOFYOUTH
		% -----
			\label{GM-LIFE-ENERGY-0001}
			
			Rejoice!

			And be glad! All those frustrations, all those hesitations, the nervousness, the burning, the butterflies of unease, the confusions, and the errors that manifest itself when your mind says ``Do this'' and your body and personality do not follow---all of these can be washed away...forever.

			Yes, you thought yourself ignorant about women, so inexperienced, so confused, and you just seemed to be following an evil circle. Maybe you were like me, a super mega dork who slept through life until now. And, perhaps like I did, you read and reread all material on women and life you could get your hands on.

			But what if I told you that you were always a Don Juan at one point in time? And I do not mean in a specific situation or hour, I mean at one point in your life you KNEW EVERYTHING ABOUT WOMEN and handled them WITH PERFECT EASE, so perfect that it was 100\% NATURAL and if anyone told you to view a website or books for `information' you would laugh yourself silly!

			Wouldn't that change everything? Wouldn't you say to yourself, ``Well! If I was the PERFECT Don Juan then I would try to remember how I THOUGHT and what I DID when I WAS that Don Juan rather than mine endless posts and books. After all, if I WAS a Don Juan, then I require only \textit{recollection} NOT \textit{revelation}.''
			
			Go find a picture of yourself when you were a young kid (say at age 6). Look at him! He is smiling gleefully without a care in the world. He doesn't know he is going to turn into the sad adult that you are now. Hormonally, the only difference between you and your youthful shadow is that you are flooded in testosterone and in a state of chemical madness. Your youthful shadow knows better of the joys of life. It is no wonder parents find their offspring such a wonder as they reframe the dull life of bills, appointments, and responsibilities with the fire of youth.

			Yes, you say that you are an \textit{adult} now. You have \textit{responsibilities} such as bills and chores and work to do. You have no time for such nonsense. Or do you?
			
			I bet when you were a kid, you were a \textit{natural} Don Juan. I bet you got all the girls in the sandbox. As a child, you knew how to treat women better than even though you are now an adult. Some of these factors include that you...
			
			-Knew girls and guys were different.

			\textit{Trendy intellectuals have this problem today! You KNEW there was MALE and FEMALE.}

			-Realized girls had cooties and could destroy a guy.

			\textit{This isn't too far off as girls can totally devour and destroy a man and his life.}

			-Knew that it was \textit{improper} to be girlish. You would (and ought to) get beaten up on the playground.

			\textit{Older people get an error in the brain called PHILOSOPHY that speaks bubble swelled words like `relativity', `revolution', freedom' and pop with the scent of rotten eggs of pious moralizations, bumper sticker arguments, and rambling dissertations. You knew, at an early age that YOU CANNOT FREE YOURSELF FROM GENDER.}

			-Girls were not to be taken seriously. After all, they are girls.

			\textit{It is the nice guy that takes the woman seriously in every and all things AND CANNOT SAY NO TO HER!}

			-Girls were to be guided, teased, because, after all, they were girls and, as such, tended to mess things up. You pulled their hair, made fun of their clothes...

			\textit{...not because an internet guy named DeAngelo told you to do so, but because you knew instinctively that it was right and proper for you to do.}

			-As a boy, you would never leave the plans up to the woman. Oh, that would be awful! As a kid, you had to say WHAT you two were going to do, WHEN, WHERE, and sometimes WHY. You had to be direct.
			
			\textit{``Why do I need a plan, Pook? Why can't we do what SHE wants to do?'' Because she has no idea what she wants to do. Example:}
			
			\textit{``When do you want to come over?''}
			
			\textit{``I don't know.''}
			
			\textit{``On Tuesday or Thursday?''}
			
			\textit{``I just don't know.''}
			
			\textit{``How about Wednesday? Is that OK for you?''}
			
			\textit{``Maybe...''}
			
			\textit{``What about Tuesday or Friday?''}
			
			\textit{``I dunno!''}
			
			\textit{*aggravated* ``Argh! I am coming to pick you up at 7:45 PM on Wednesday.''}
			
			\textit{*sweetly* ``OK!''}
			
			-You did not get into serious talks with a girl. You did not turn her into Oprah. You did not try to impress her with how `intellectual' you are. You probably hit her, cried ``Tag!'', ran off, and she would chase after you. You would get on the swings. You would push her off the slides. On the see-saw, you would try to fall as fast as possible to catapult her away.
			
			\textit{Everything you did with a girl was ACTION dates. Nice Guys try to cook the lady dinner. YOU wouldn't even DARE do such a thing when you were young. Now I know why some of my best dates are ones as simple as taking the girl to the park and run around like little kids. Action! Action! Action!}

			-As a kid you loved to sing. You loved to laugh.

			\textit{What do you do now? You are so uptight that you wouldn't catch yourself dead singing outside your home. And what happened to that happy laughter that marked your childhood? Why are you taking everything so seriously now!?}
			
			-You had toys and loved to play with them. When the girls entered your sphere, you insisted on playing with YOUR toys. You would ride your bike at death defying speeds. You played with cranes and Tonka dump trucks. You LOVED firetrucks.

			\textit{When you grow up, your Hot Wheels cars turn into super fast sport cars (which you still drive at death defying speeds). Your cranes and dump trucks turn into the big ones and you still love the firetrucks. Compare this to the Nice Guy who does not understand the beauty of construction or get scared at the idea of firetrucks. There is a reason why women LOVE firemen.}

			-You embraced your imagination. If you played with a girl, it was to be \textit{on your terms}. You will not find a boy that says to a girl, ``Whatever YOU want to do.'' I remember flicking caterpillars on girls (and they loved it!). You would point to the girl and go, ``You are the detective and I am the cop.'' And then you pointed to others and go, ``Look at these villains! Come, we must round them up!'' And the girls joyfully played the part.

			\textit{What do you do NOW? ``Let's go get dinner.'' BORING.}

			-You had one eternal enemy in childhood, boredom. Like a void, it encroached on you in school, ensnared you with a stupid trip with the parents, and enveloped you as you stepped on the school bus. You embraced every chance for play.

			\textit{If there is ONE thing a man must NEVER do to a woman, it is this: DO NOT BORE HER. Make her happy, make her angry, make her laugh, ANYTHING but bore her. That enemy, boredom, is back and women are looking for you to strike it down. Embrace your youth and live again for the first time.}
			
			-Even at your young age, you were aware of fashion. Your mother was perplexed at why you couldn't wear THAT shirt or put on THOSE shoes. You knew the importance clothes and appearance had.

			\textit{Nice Guys and chumps try to say nothing for appearance and say, ``she will like me for who I am'', and, with the same breath, ignore the chicks who do nothing appearance wise (big whale chicks, pimply chicks, and such) while getting shot down at the real women.}

			-When you were young, your father was a DEMI-GOD. You both feared and loved him. He could be playful when he wanted to. But, always, he was a SOURCE of STRENGTH, always confidant, and always seemed to know the solution to any problem you came across. This feeling of awe you thought of your father is the perfect definition of a MAN.

			\textit{When you strive to be that same towering figure, which seemed to have solutions to all problems, confidence for all troubles, know how for all messes, stability for all storms, women will react to you in the same way. It is said, ``Women want to marry their fathers.'' But this phrase has been taken completely out of context. Women want to marry THAT guy, that MAN they knew when they were a little girl. You can only understand and become that man through the eyes of the young boy.}
			
			Now some women will protest this advocacy of looking at women as little girls. This works especially well with the YOUNGER chicks (which is what most guys here want). This also solves a legal mystery: why were women from 1800s and down treated, BY LAW, as children? We know the answer now). (And I would say to women to treat men the same way. Women are most charming when they view us as boys.)

			``He, whom the gods love, grows young,'' said the ancient Greeks. Look at that picture of the young you! Now look in the mirror. The blazing light in the youth's eyes, the curiosity, and the wonders he saw at Nature and life, the joy he loved at any occasion! Are they still part of you? Look into the mirror. Is the same light in your eyes?

			If not, then you know what you need to do. You have come to the solution of your Don Juan troubles. Legends spoke of a Fountain of Youth that turns the old young. Conquistadors prowled continents hunting for this magical facet of Nature. This legend of the Fountain of Youth was like a psychological splinter in the minds of so many men, driving them across the world through dangers and storms to obtain such a treasure. But it was not in the world, it was in us. Who knew women are the key that seems to unlock all of Nature's mysteries?

			Doesn't this clear the stormy air of confusion? (especially with younger `immature' women!) When you are with a woman and a thousand Don Juan philosophies and tactics come to thwart your peace, remember the kid that you were... and how he looked on life. Yes, she might be a twenty-something vixen. But underneath those milk-sacs and fat deposits that drive your chemicals mad, is a little girl. Rather than being nervous about some date, view the date as if you were seven years old. Everything becomes simple and fun (as it should be!).

			One thing is for certain, women go NUTS over a guy who keeps his boyhood charm as women want an ESCAPE. They do not want to hear your views on the world. They want to have the happiness and fun of their childhood back (as everyone does!). They will FIGHT for the rare men who truly live like this (this also explains the mystery why uneducated men often seem to do BETTER with women than many men with PhDs who are so `smart' they intellectualized life out of existence).

			As we know, cell division error, aided by free radicals, accumulates errors throughout the body as time passes. A person of 80 obviously has more errors than that of 40. His tissues start to fail, making his organs fail, then entire organ systems collapse, until life is snuffed out.

			In the same way, our minds start out pristine and pure (that of the child!) and everything is playful, fun, and simple. But as time passes, philosophies and bitter memories accumulate. The diseased person looks at life only through the philosophical lens, living a life of past memories, and so his life decays and decays until there is no life within him.

			\textit{Pook hands you the chalice.} \hl{Drink and pass the cup around. This water from the Fountain of Youth will wash away these errors, all those `frames' your bitter memories eat at your mind. Drink to the girls that rejected you viciously. Drink to the `macho' guys that beat you up in school. Drink to the chains of routines, errands, and appointments. Drink to your ambition, drink to your melancholy, drink to your loneliness. Drink to your heart's content and pass the cup around. Your mind has now absorbed the blessed waters of the Fountain of Youth. \textit{You are now the light of the world!} Your life is now young, fruitful, fun, easy, simple, and your body will reflect it as well.}

			\textbf{The world is now your sandbox.} Rejoice! Many people get swallowed up in their vanity, all believing themselves `brilliant' and `smart'. They follow their philosophies unquestioningly and their lives walk on the quicksand of melancholy. But... they see the new you and it smashes all their philosophies to bits. They ask each other, ``How can he be so happy? How can he be so ALIVE? Let us study him and we will write articles and manifestos on him.'' But you know the truth. Their ambitions have consumed them while you look on life with youthful eyes...

			...and live in a Child's Paradise.
			
			\begin{center}
			* * * 
			\end{center}
			
			Quote:
			
				\begin{quote}
				Originally posted by ShortTimer
				
				\textit{Indeed, maybe the childhood Pook speaks of is true of him but it seems more like some idealistic version of the past that never really existed. I too recall my childhood with torment, if I could I wouldn't live it again. Maybe this metaphor will work for some people tho...}
				\end{quote}
				
			\noindent You guys are thinking too much of this.

			There is no \textbf{Time} in Sexuality. In the realm of sexuality uncorrupted with `intellect' and `politics' and all, which is Womaniverse, time simply doesn't exist.

			This is not looking at your childhood as if it were some `Golden Age'. To the contrary, your entire life is the Golden Age. The problem isn't because you aged (because in sexuality, time doesn't exist), the problem with many is they've become so super-intellectual or so `educated' that it corrupts their sexuality.

			Your body doesn't get corrupted by Time. It gets corrupted by cell division error and free radicals. In the same manner, your mind/soul doesn't age. It gets corrupted with `intellect', `philospphies', and `education'.

			I always wondered WHY girls could be attracted to hard bodies, happy smiles, but NEVER the genius. Now we know why.

			Years ago, I was at my `most natural state' when I would rant about philosophy, politics, economics, and so on.

			As you can imagine, this bored the women (women want smart guys but not `melancholy' filled philosophic guys). My actions were so `intellectualized' that women saw through everything (or used me in some way).

			This is an anti-intellectual post (very tame of what will follow). Many people sense that Time is going faster and faster and faster for them, as if Life is speeding up. Things become more frantic, responsibilities become more stacked, and the geometry of rules and regulations grow narrower and narrower around your life until nothing remains.

			The reason why Time is speeding up is because your life's relationship to sexuality is going down.

			This isn't about bringing your childhood into the present. It is to CLEANSE your mind and soul of all this philosophic DUST and bitter memory GARBAGE that has accumulated.

			We cannot reverse time with our bodies (thanks to the cell-division error). But we can drink from this fountain of youth for our own minds.

			With one sip, the `complexity' of women just vanish!

			With two sips, you no longer want to `talk' and `intellectualize' with girls, you just want to run around and have fun.

			With three sips, you look at all these `manifestos' of women and how to get women with disbelief.

			This Fountain of Youth is NOT immaturity. It is NOT mediocrity.

			I find it amazing how people go through so much trouble to make sure their body isn't contaminated by disease, germs, and the like. Yet, there is no such consideration for their mind and soul with politics, philosophies, and perversions of all sort to infect and tear down your sexuality, and thus, your life.

			Think about it. Do you want to have a relationship based on talking about DNA?

			The childhood example applies because \textit{no child} intellectualizes life. It was only years and years later that I decided I was `smart' and so read and read and got filled with a lot of BS that life slipped past me.

			Most people's sexuality deteriorates. Some people's deteriorates more than others. But when one of these people see a Don Juan, oh, they start intellectualizing what is `really' going on. So all these `theories' are born which is just the intellectual's excuse to not face the truth about himself….

			\textit{That HE was not so `brilliant' after all!}
		
		% -----
		\subsubsection{0002 | Walter White after his cancer diagnosis}
		% -----
			\label{GM-LIFE-ENERGY-0002}
			
			I am awake...we start tomorrow. (\textit{Breaking Bad}, episode 1, season 1)
			
	% ----------------------
	\subsection{Time and its use} % *TIME
	% ----------------------
		\label{GM-LIFE-TIME}
			
		% -----
		\subsubsection{0001 | Burkeman on time speeding up, and how to slow it down}
		% -----
			\label{GM-LIFE-TIME-0001}
			
			Meanwhile, no catalog of our time-related troubles would be complete without mentioning that alarming phenomenon, familiar to anyone older than about thirty, whereby time seems to speed up as you age---steadily accelerating until, to judge from the reports of people in their seventies and eighties, months begin to flash by in what feels like minutes. It's hard to imagine a crueler arrangement: not only are our four thousand weeks constantly running out, but the fewer of them we have left, the faster we seem to lose them (4TW PDF 15).
				\begin{compactitem}
					\item See also Easter on \hyperref[GM-LIFE-TIME-0006]{learning and new experiences} and how it affects our perception of time.
				\end{compactitem}
			
		% -----
		\subsubsection{0002 | Burkeman on the curse of busyness}
		% -----
			\label{GM-LIFE-TIME-0002}
			
			The problem isn't exactly that these [time-management] techniques and products don't work. It's that they do work---in the sense that you'll get more done, race to more meetings, ferry your kids to more after-school activities, generate more profit for your employer---and yet, paradoxically, you only feel busier, more anxious, and somehow emptier as a result. (4TW PDF 16)
			
			Instead, life accelerates, and everyone grows more impatient. It's somehow vastly more aggravating to wait two minutes for the microwave than two hours for the oven ---or ten seconds for a slow-loading web page versus three days to receive the same information by mail. (4TW PDF 17)
			
			I did get better at racing through my to-do list, only to find that greater volumes of work magically started to appear. (Actually, it's not magic; it's simple psychology, plus capitalism). (4TW PDF 18)
			
			But now here we get to the heart of things, to a feeling that goes deeper, and that's harder to put into words: the sense that despite all this activity, even the relatively privileged among us rarely get around to doing the right things. We sense that there are important and fulfilling ways we could be spending our time, even if we can't say exactly what they are---yet we systematically spend our days doing other things instead. (4TW PDF 19) 

			Our days are spent trying to ``get through'' tasks, in order to get them ``out of the way,'' with the result that we live mentally in the future, waiting for when we'll finally get around to what really matters---and worrying, in the meantime, that we don't measure up, that we might lack the drive or stamina to keep pace with the speed at which life now seems to move. (4TW PDF 19--20) 
			
			The real problem---or so I hope to convince you---is that we've unwittingly inherited, and feel pressured to live by, a troublesome set of ideas about how to use our limited time, all of which are pretty much guaranteed to make things worse. (4TW PDF 23)
			
		% -----
		\subsubsection{0003 | Burkeman on how you'll never feel free enough to do what you want}
		% -----
			\label{GM-LIFE-TIME-0003}
			
			Productivity is a trap. Becoming more efficient just makes you more rushed, and trying to clear the decks simply makes them fill up again faster. Nobody in the history of humanity has ever achieved ``work-life balance,'' whatever that might be, and you certainly won't get there by copying the ``six things successful people do before 7:00 a.m.'' The day will never arrive when you finally have everything under control---when the flood of emails has been contained; when your to-do lists have stopped getting longer; when you're meeting all your obligations at work and in your home life; when nobody's angry with you for missing a deadline or dropping the ball; and when the fully optimized person you've become can turn, at long last, to the things life is really supposed to be about. (4TW PDF 20--21) 
			
		% -----
		\subsubsection{0004 | Deida on how you'll never feel free enough to do what you want}
		% -----
			\label{GM-LIFE-TIME-0004}
			
			Most men make the error of thinking that one day it will be done. They think, ``If I can work enough, then one day I could rest.'' Or, ``One day my woman will understand something and then she will stop complaining.'' Or, ``I'm only doing this now so that one day I can do what I really want with my life.'' The masculine error is to think that eventually things will be different in some fundamental way. They won't. It never ends. As long as life continues, the creative challenge is to tussle, play, and make love with the present moment while giving your unique gift. 

			It's never going to be over, so stop waiting for the good stuff. As of now, spend a minimum of one hour a day doing whatever you are waiting to do until your finances are more secure, or until the children have grown and left home, or until you have finished your obligations and you feel free to do what you really want to do. Don't wait any longer. Don't believe in the myth of ``one day when everything will be different.'' Do what you love to do, what you are waiting to do, what you've been born to do, now. 

			Spend at least one hour a day doing whatever you simply love to do---what you deeply feel you need to do, in your heart---in spite of the daily duties that seem to constrain you. However, be forewarned: you may discover that you don't, or can't, do it; that, in fact, your fantasy of your future life is simply a fantasy. 

			Most postponements are excuses for a lack of creative discipline. Limited money and family obligations have never stopped a man who really wanted to do something, although they provide excuses for a man who is not really up to the creative challenge in the first place. Find out today whether you are willing to do what it takes to give your gift fully. As a first step, spend at least an hour today giving your fullest gift, whatever that is for today, so that when you go to sleep at night you know you couldn't have lived your day with more courage, creativity, and giving...

			The world and your woman will always present you with unforeseen challenges. You are either living fully, giving your gift in the midst of those challenges, even today, or you are waiting for an imaginary future which will never come. Men who have lived significant lives are men who never waited: not for money, security, ease, or women. Feel what you want to give most as a gift, to your woman and to the world, and do what you can to give it today. Every moment waited is a moment wasted, and each wasted moment degrades your clarity of purpose. (WSM PDF 23--24)
			
		% -----
		\subsubsection{0005 | Burkeman on comparing and keeping up with others}
		% -----
			\label{GM-LIFE-TIME-0005}
			
			It turns out that when people make enough money to meet their needs, they just find new things to need and new lifestyles to aspire to; they never quite manage to keep up with the Joneses, because whenever they're in danger of getting close, they nominate new and better Joneses with whom to try to keep up. As a result, they work harder and harder, and soon busyness becomes an emblem of prestige. Which is clearly completely absurd: for almost the whole of history, the entire point of being rich was not having to work so much. Moreover, the busyness of the better-off is contagious, because one extremely effective way to make more money, for those at the top of the tree, is to cut costs and make efficiency improvements in their companies and industries. That means greater insecurity for those lower down, who are then obliged to work harder just to get by. (4TW PDF 18)
			
		% -----
		\subsubsection{0006 | Easter on learning new things and being in the moment}
		% -----
			\label{GM-LIFE-TIME-0006}
			
			Learning new skills---particularly the ones humans needed for millions of years that require us to use our mind and body---would stay with me beyond Alaska in a very Zen, the-path-is-the-goal type of way. There was all the new expertise I was picking up. But learning new skills is also one of the best ways to enhance awareness of the present moment, with no burning incense, Buddhist mantras, or meditation apps involved.
			
			My months of preparation changed much of that. New situations kill themental clutter. In newness we're forced into presence and focus. This is because we can't anticipate what to expect and how to respond, breaking the trance that leads to life in fast forward. Newness can even slow down our sense of time. This explains why time seemed slower when we were kids. Everything was new then and we were constantly learning.

			Psychologist William James wrote about this in his 1890 work, \textit{ThePrinciples of Psychology}: ``The same space of time seems shorter as we grow older...In youth we may have an absolutely new experience, subjective or objective, every hour of the day. Apprehension is vivid, retentiveness strong, and our recollections of that time, like those of a time spent in rapid and interesting travel, are of something intricate, multitudinous, and long-drawn-out. But as each passing year converts some of this experience into automatic routine that we hardly note at all, the days and the weeks smooth themselves out in recollection to contentless units, and the years grow hollow and collapse.'' (TCC PDF 65)
		
	% ----------------------
	\subsection{Creativity} % *CREATIVITY
	% ----------------------
		\label{GM-LIFE-CREATIVITY}
		
		% -----
		\subsubsection{0001 | Easter on the importance of boredom for creativity}
		% -----
			\label{GM-LIFE-CREATIVITY-0001}
			
			There's another massive benefit to boredom beyond making us more psychologically robust and resilient. Finding a different outlet for boredom also lets us tap into creativity. 

			In an interview with Bill Simmons, the acclaimed and prolific screenwriter Aaron Sorkin summed up this phenomenon when he talked about the first time he ever wrote for fun: ``It was one of those nights in New York where it feels like everyone has been invited to a party you weren't invited to. I didn't have three dollars in my pocket. In [my] apartment was a semiautomatic typewriter. Electric keys and a manual return. The TV was broken, the stereo was broken. The only thing to do was to put a piece of paper in that typewriter and to start typing. Pure boredom. It was the first time I wrote for fun…and I loved it. I stayed up all night writing and I feel like that night has never ended.'' 

			Sorkin's takeaway is that we should learn to deal with boredom, and then discover ways to overcome it that are more productive and creative than watching a YouTube video or scrolling through Instagram. 

			Research from the 1950s backs up Sorkin's connection between boredom and creativity. One team of UK researchers had people do something astoundingly boring: read a phone book for 15 minutes. The bored people then took a standardized creativity test, like coming up with odd uses for a Styrofoam cup, and the Remote Associates Test (RAT), where three words are thrown out and we have to figure out their common denominator (e.g., call + pay + line = phone; motion + poke + down = slow). Compared to a nonbored group, the bored people gave significantly more answers on both tests, and the answers were also considerably more creative. Other studies have found the same phenomenon. (Except in those studies the researchers bored people by having them watch a screensaver or separate a pile of beans by color.) 

			``But now people want to say that boredom makes you more creative,'' said Danckert. ``I call bullshit on that. Boredom doesn't make you more creative. It just tells you ‘do something!' '' And when that ``something'' is letting our mind revive unfocused mode---or sitting down to write a screenplay---rather than blanketing it with the exact same media that everyone else is consuming, we begin to think, quite literally, on a different wavelength. That's what creativity requires. 

			And now we've killed off one of the main drivers of creativity: mind- wandering...The scientist blames our hurried, overscheduled lives and ``ever increasing amounts of [time] interacting with electronic entertainment devices.'' 

			My mind feels like it's riding a different wavelength than the one it typically does back home. It's more a ripple than a riptide. Despite the cold, wind, and rough ground, my stress levels are nonexistent. Interesting new ideas are bubbling out of the ether. And I'm oddly appreciative---of the world around me but also of things back home. Like my wife...I believe all these feelings have something to do with allowing my mind a moment of rest. Maybe when I get home, instead of thinking the oft- repeated ``less phone,'' it might be more productive to think ``more boredom.'' (TCC PDF 103--108)
			
	% ----------------------
	\subsection{Growth and discomfort} % *GROWTH
	% ----------------------
		\label{GM-LIFE-GROWTH}
		
		% -----
		\subsubsection{0001 | Agnes Callard on constantly learning}
		% -----
			\label{GM-LIFE-GROWTH-0001}
			
			The things people long for are: safety and security; fancy vacations and luxury goods; honor, power and acclaim; the warmth of family life and human connection. They want these things even when they don't have them---often, the less they have them, the more they want them. People don't long for intellectual goods. You know the joys of intellectual engagement by experiencing them, and as you step away from them they fade from view. There are strange people who somehow, through a series of accidents, get and stay keyed onto intellectual goods on their own---the autodidacts I mentioned earlier---but the rest of us need constant help reorienting, because just about every worldly temptation pulls us in the opposite direction. (``The Real College Scandal'')
		
		% -----
		\subsubsection{0002 | Easter on comforts we have}
		% -----
			\label{GM-LIFE-GROWTH-0002}
			
			I wasn't a completely new person and I'd never be confused with Mr. Rogers. But I was more aware, which allowed me to see that I was still surrounded in comfort. I was marinating in the stuff. Except that these were less acutely destructive but potentially more insidious forms of it. I just had to take a look at my everyday life. I was comfortable, quite literally, every single moment. 

			I awoke in a soft bed in a temperature-controlled home. I commuted to work in a pickup with all the conveniences of a luxury sedan. I killed any semblance of boredom with my smartphone. I sat in an ergonomic desk chair staring at a screen all day, working with my mind and not my body. When I arrived home from work, I filled my face with no-effort, highly caloric foods that came from Lord knows where. Then I plopped down on my overstuffed sofa to binge on television streamed down from outer space. I rarely, if ever, felt the sensation of discomfort. The most physically uncomfortable thing I did, exercise, was executed inside an air-conditioned building as I watched cable news channels that are increasingly bent on confirming my worldview rather than challenging it. I wouldn't run outside unless the conditions were, well, comfortable. Neither too hot, too cold, nor too wet. 

			What could cleansing myself of all these other comforts do for me? (TCC PDF 19)
			
			Except that our original comforts were negligible and short-lived, at best. In an uncomfortable world, consistently seeking a sliver of comfort helped us stay alive. Our common problem today is that our environment has changed, but our wiring hasn't. And this wiring is deeply ingrained. (TCC PDF 20)
			
			Finally, on June 29, 2007, boredom was pronounced dead, thanks to the iPhone. And so our imaginations and deep social connections went with it. (TCC PDF 23)
			
			For most modern Americans, ``stress'' is so often ``This traffic is going to make me late to my yoga class'' stress. Or ``Is my neighbor making more money than me?'' stress. Or ``This spreadsheet is going to take forever'' stress. Or ``If my child doesn't get into an Ivy League school we will all live lives of complete and utter nothingness'' stress. It's first-world stress. (TCC PDF 24)
			
			So, yes, we don't have to deal with discomforts like working for our food, moving hard and heavy each day, feeling deep hunger, and being exposed to the elements. But we do have to deal with the side effects of our comfort: long-term physical and mental health problems. We lack physical struggles, like having to work hard for our livelihoods. We have too many ways to numb out, like comfort food, cigarettes, alcohol, pills, smartphones, and TV. We're detached from the things that make us feel happy and alive, like connection, being in the natural world, effort, and perseverance. 

			We seem to know something is amiss. One poll found just 6 percent of Americans believe the world is improving. Some anthropologists, in fact, argue that humans were happier in all the time leading up to about 13,000 years ago. People then had simpler needs that were easier to fulfill and were more able to live in the present. 

			Comforts and conveniences are great. But they haven't always moved the ball downfield in our most important metric: happy, healthful years. Perhaps existing only in our increasingly overly comfortable, overbuilt environment and always obeying our comfort drives has had unintended consequences and caused us to miss profound human experiences. There are conditions that humans evolved to live in and experiences we were meant to have that are no longer germane to our lives. This has undoubtedly changed us, often not for the best. (TCC PDF 25)
			
			It explains that as we experience fewer problems, we don't become more satisfied. We just lower our threshold for what we consider a problem. We end up with the same number of troubles. Except our new problems are progressively more hollow...We are always moving the goalpost. There is, quite literally, a scientific basis for first-world problems...Call it comfort creep. When a new comfort is introduced, we adapt to it and our old comforts become unacceptable. Today's comfort is tomorrow's discomfort. This leads to a new level of what's considered comfortable. (TCC PDF 27, 28)
				\begin{compactitem}
					\item See \hyperref[GM-LIFE-ENERGY-0005]{Burkeman} on comparing and keeping up with others, who makes a very similar point.
				\end{compactitem}
				
		% -----
		\subsubsection{0003 | Easter on the necessity of pushing yourself to live at your edge, and misogis}
		% -----
			\label{GM-LIFE-GROWTH-0003}
			
			``In modern society, however,'' Elliott said, ``it's suddenly possible to survive without being challenged. You'll still have plenty of food. You'll have a comfortable home. A good job to show up to, and some people who love you. And that seems like an OK life, right? ``But,'' he said, sweeping his arm to create a big imaginary circle that encompassed the trail and foliage flanking it. ``Let's say your potential is this big circle.'' 

			Then he pulled his hands into his chest and made a dinner-plate-size circle in the exact middle of the much larger circle. ``Well, most of us live in this small space right here. We have no idea what exists on the edges of our potential. And by not having any idea what it's like out on the edge...man, we really miss something vital.'' 

			A salty wind was blowing over the ocean and sidling up into the hills. It passed over my sweat-soaked T-shirt as Elliott continued. ``I believe people have innate evolutionary machinery that gets triggered when they go out and do really fucking hard things. When they explore those edges of their comfort zone.'' 

			``Misogis are that same concept. Except for the modern condition. In misogi we're using the artificial, contrived concept of going out and doing a hard task to mimic these challenges that humans used to face all the time. These challenges that our environment used to naturally show us that we're so removed from now,'' he said. ``Then when we return to the Wild West of our everyday lives we are better for it. We have the right tools for the job.'' 

			``Misogi is not about physical accomplishment,'' said Parrish. ``It asks, `What are you mentally and spiritually willing to put yourself through to be a better human?' Misogis have allowed me to let go of fear and anxiousness, and you can see that in my work.'' (TCC PDF 39--40)
			
		% -----
		\subsubsection{0004 | Easter in getting yourself uncomfortable}
		% -----
			\label{GM-LIFE-GROWTH-0004}
			
			But all of that silent, solitary time running, riding, or swimming--- becoming comfortable with discomfort, persisting despite all of his biological impulses telling him to slow down or tap out---had remodeled his psyche. ``Endurance sports gave me some understanding of what it was to push to deeper levels and find new layers within myself,'' he told me. ``When I stopped doing triathlons, I still had this sense of adventure. This need to explore those edges where I'd find a new, better part of myself.'' 

			``It was all just to see if I could. I'd get to what I thought was my edge, but I'd keep going. Then eventually I'd realize I was way past my old edge and still going. And so that edge was now in a different place than when I started. And that was so satisfying, so satisfying.'' 

			The misogis have continued---one a year---and Elliott credits them with his ability to affect things in his personal and professional life. ``Misogis can show you that you had this latent potential you didn't realize, and that you can go further than you ever believed. When you put yourself in a challenging environment where you have a good chance of failing, lots of fears fade and things start moving.'' (TCC PDF 42, 44)
			
		% -----
		\subsubsection{0005 | Easter on designing challenges to push yourself to discomfort: a rite of passage}
		% -----
			\label{GM-LIFE-GROWTH-0005}
			
			``We're generally guided by the idea that you should have a fifty percent chance of success---if you do everything right,'' he said. ``So if you decided you wanted to run a twenty-five-mile trail, and you're preparing by working up to a twenty-mile training run and doing thirty-five or forty miles a week of running...that's not a misogi. Your chance of failure is too low. But if you've never run more than ten miles, think you could probably run fifteen, but are iffy on whether you could run twenty...then that twenty-five miles is probably a misogi.''

			Engaging in an environment where there's a high probability of failure, even if you execute perfectly, has huge ramifications for helping you lose a fear of failing. Huge ramifications for showing you what your potential is.

			Men and women in these cultures [hunter-gatherer, older] undertook a physical, nature-based rite of passage.
			
			``The idea of a rite of passage is that the elders are seeing in you the potential to rise up and achieve this really important, challenging thing that is going to benefit you and everyone around you on many levels,'' said Elliott. ``They're saying, `We think you're ready, but you're really going to have to dig deep and find your shit.' ''

			He found that these processes---whether walking around the outback, hunting a lion in Kenya, tripping out on the Columbia River Plateau, or perhaps, even, undergoing a misogi---all have three key elements. The first is separation. The person exits the society in which they live  and ventures into the wild. The second is transition. The person enters a challenging middle ground, where they battle with nature and their mind telling them to quit. The third is incorporation. The person completes the challenge and reenters their normal life an improved person. It's an exploration and expansion of the edge of a person's comfort zone.
				\begin{compactitem}
					\item These are the rules for the misogi.
				\end{compactitem}
			
			``Yes. And the reason for this is because the more quirky the misogi, the less chance you can compare it to anything else,'' he said. ``It's important to take on challenges that are your challenges. Misogi is you against you. It's against this phenomenon of `Oh, that guy did this thing in this amount of time and I'm going to try to do it faster.' Because that's comparison shopping. And that's just such a shitty way to go through life.'' Parrish talked at length about this guideline. He summed it up like this: ``When you remove superficial metrics you can accomplish way more.'' Which brought Elliott to guideline two: Don't advertise misogi. It's OK to talk about misogi with friends and family. But you don't Tweet, Instagram, Facebook, or boast about misogi. 

			``Everyone today has such outward-facing lives,'' said Elliott. ``They do stuff so they can post on social media about some badass thing they did to get a bunch of likes. 

			``Misogis are inward facing,'' he said. ``A big part of the value proposition is that I'm going to do something that's really uncomfortable. I'm going to want to quit. And it's going to be hard not to quit because no one is watching. But I'm not going to quit because I'm watching. And then I can reflect back on how I was the only person watching myself and I still rose to the occasion in a big way. There's some deep satisfaction in that. Did you really do what you think is the right thing when you were the only person watching? Or do you need an audience or a big pat on the back for that? Are you not important enough to do it for you? We had this guideline before social media, and it seems more relevant today.''
			
			Finding discomfort in Las Vegas was easy. All I had to do was walk into the desert. Hiking and running in the summer felt like exercising in a furnace. But these hikes and runs were also an oasis away from my daily life in the built environment. I'd toss on my gargantuan \$400 boots (they needed breaking in, after all) and trudge through red rock canyons, black rock desert, Joshua tree forests, or piney highlands for a few hours each weekend. These natural environments acted like a pressure washer on my mind, clearing out the week's grime. Who needs to chat up a \$100/hour therapist when there are long, quiet, empty trails waiting to be walked? (TCC 45--59, 63)
			
		% -----
		\subsubsection{0006 | Easter on the importance of being ok with silence}
		% -----
			\label{GM-LIFE-GROWTH-0006}
			
			We have in fact become so used to living with noise that most of us now find comfort in constant blare, according to a scientist in Australia. The researcher had hundreds of students spend a little time in the quiet and write about their experiences. Nearly every student said the silence made them uneasy. ``The lack of noise made me uncomfortable, it actually seemed foreboding,'' wrote one student. ``Perhaps because media consistently surrounds us today, we have a fear of peace and quiet,'' wrote another. 

			Another survey found that Americans increasingly see the TV not as an entertainment device but as a companion. More than half of us keep the TV on while we work, cook, and do chores because we feel uncomfortable in silence. Silence-induced discomfort is a new, learned behavior, those Australian scientists think.
			
			Hearing the natural sounds we evolved in seems to strike a calming note within us. Scientists in the UK, for example, found that people who listened to nature sounds like water and wind reduced their stress levels significantly more than those who listened to artificial noises. (TCC PDF 126, 130)
			
	% ----------------------
	\subsection{The world and distractions} % *WORLD
	% ----------------------
		\label{GM-LIFE-WORLD}
		
		% -----
		\subsubsection{0001 | Easter on the landscape of despair}
		% -----
			\label{GM-LIFE-WORLD-0001}
			
			Some mental health researchers today call our concrete, sprawling environments ``landscapes of despair.'' But the Industrial Revolution spurred a great migration into cities with the promise of secure jobs. We haven't turned back since. Yet, interestingly enough, money doesn't seem to overcome the rural/urban happiness gap. Studies show that even dirt-poor people who live in rural China report being happier than infinitely wealthier Chinese city-dwellers. (TCC PDF 69)
				\begin{compactitem}
					\item Note that the internet is also a landscape of despair---it isn't just a physical place.
				\end{compactitem}
				
			I can't help but think that in today's increasingly hyperconnected and tribal society---where we define ourselves by the group or movement we belong to---it's not a bad idea to occasionally be alone. Removed from anyone. I'm talking about time with yourself, unidentified with anything. The Buddha, Lao-Tzu, Moses, Milton, Emerson, and many more have spoken highly of the benefits of solitude.

			A growing field of scientists today think that these solitude-seekers were onto something. Building ``the capacity to be alone'' may be just as important for you as forging good relationships. ``The capacity to be alone is essentially the ability to be alone with yourself and not feel uncomfortable or like you have to distract yourself,'' said Matthew Bowker, PhD, a professor of psychology at Medaille College. (TCC PDF 78--79)
			
		% -----
		\subsubsection{0002 | Easter on the benefits of solitude}
		% -----
			\label{GM-LIFE-WORLD-0002}
			
			``Building the capacity to be alone probably makes your interactions with others richer. Because you're bringing to the relationship a person who's actually got stuff going on in the inside and isn't just a connector circuit that only thrives off of others.'' 

			Research backs solitude's healthy properties. It's been shown to improve productivity, creativity, empathy, and happiness, and decrease self- consciousness. 

			``Social connection is obviously critical,'' said Bowker. ``But it can be dangerous if your social connections ever go away and you don't have yourself to fall back on. If you develop that capacity to be alone, then instead of feeling lonely, you could see solitude as an opportunity to have a meaningful and enjoyable time to get to know yourself a little better. To essentially build a relationship with yourself. I know this sounds cheesy, but it's critical. I think a goal we should all have is to try to transform feelings of loneliness into feelings of rich solitude.'' (TCC PDF 80)
			
		% -----
		\subsubsection{0003 | Easter on the importance of being bored}
		% -----
			\label{GM-LIFE-WORLD-0003}
			
			The 11 hours and 6 minutes of attention we're handing over to digital media isn't free. It's all spent in focused mode. Think of this focused state like lifting a weight, and the unfocused state like resting. When we kill boredom by burying our minds in a phone, TV, or computer, our brain is putting forth a shocking amount of effort. Like trying to do rep after rep after rep of an exercise, our attention eventually tires when we overwork it. Modern life overworks the hell out of our brains.

			Our collective lack of boredom may be causing us to reach near-crisis levels of mental fatigue. Research shows that the onslaught of screen-based media has created Americans who are ``increasingly picky, impatient, distracted, and demanding,'' as one media analyst put it. These terms fall under the umbrella of ``insufferable.'' And overworked, undermaintained minds are linked to depression, life dissatisfaction, the perception that life goes by quicker, and increasingly missing the beauty of life that only presents itself when we allow our mind to wander and be aware of something other than a screen.
			
			The way we dealt with boredom before we began surrounding ourselves in constant comfort delivered benefits that are essential for our brain health, productivity, personal sanity, and sense of meaning. But there's been a cosmic shift in boredom. The way we now deal with it is ``like junk food for your mind,'' said Danckert.
			
			These Silicon Valley savants have big data telling them exactly what tricks will grab us, and morons like me don't stand a chance. 
			
			``There's a trigger, a behavior, and a reward,'' said Brewer. ``But this brain process can get hijacked in the modern day. The trigger instead of food is boredom. And the behavior is going on YouTube or checking our news feed or Instagram. And that distracts us from the boredom. We become excited and get a hit of dopamine, which is a reward.
			
			``The paradox is that these mechanisms that helped keep us alive are now hurting our health,'' he explained. ``We have less tolerance for distress. If we feel something unpleasant, like boredom, typically we would have to just be with that unpleasantness, and then we'd find a productive outlet. But we don't have to do that anymore. We can use our phone to distract ourselves.'' Or, as Danckert put it, we simply consume more ``junk food for the mind.'' 

			Each time we reflexively take out our phone or turn on a computer orTV to kill boredom, it attaches another tiny anchor to our stress tolerance, dragging it lower. Scientists at Oregon State University found that daily stressors like lines and waits can improve our resistance to some brain diseases if we simply suffer through them and shrug them off. More of these everyday stressors are actually better for our brain. (TCC PDF 96, 97, 101)
				\begin{compactitem}
					\item See also Easter on the \hyperref[GM-LIFE-CREATIVITY-0001]{importance of boredom for creativity}.
				\end{compactitem}
				
		% -----
		\subsubsection{0004 | Easter on the importance of being in nature}
		% -----
			\label{GM-LIFE-WORLD-0004}
				
			After a handful of techless days trudging through Alaska, I was beginning to buy in to the theory. My brain was feeling less hunkered down in its typical foxhole---a state that I'd compare to a roadrunner on crystal meth, dementedly zooming from one thing to the next---and more like it belonged to a monk after a month at a meditation retreat. I just feel…better. Wilson put my feelings this way: ``Nature holds the key to our aesthetic, intellectual, cognitive, and even spiritual satisfaction.''
			
			A walk in the woods is free. And to my knowledge it isn't linked to spastic, unplanned movements or urges to hurry down to the corner store to spend your life savings on scratch-off lottery tickets and then bang the clerk who sold them to you. Perhaps most delightful of all, being treated by nature won't require you to haggle with some stonewalling health insurance company representative...Of course, a walk in the woods only becomes mind medicine so long as the phone is away and also not beaming information into our ears. (TCC PDF 112--116)
			
	% ----------------------
	\subsection{Food and eating} % *FOOD *EATING *DIETING
	% ----------------------
		\label{GM-LIFE-FOOD}
		
		% -----
		\subsubsection{0001 | Easter on the discomfort of eating less}
		% -----
			\label{GM-LIFE-FOOD-0001}
			
			And here's the thing: Humans figured out that the amount of food weeat has something to do with our body size about 2,300 years ago. We could have stopped the research then. But we've since spent billions proving and re-proving that the ancients had it right: Eating less of something---carbs, fat, sugar, etc.---causes us to…eat less and therefore lose weight. And we also know now that all diets work until they don't. Which is, according to one large survey in the UK, an average of 5 weeks, 2 days, and 43 minutes in. That's when we give up and slowly slide back into our previous state.

			Why? This is when the discomfort sets in. People usually fail after a handful of weeks because their bodies fight to bring them back to their starting point. When your body fat drops enough, your brain responds by making you hungrier while at the same time decreasing how satisfying your meals are. A team at the NIH recently found that for every two pounds a person loses, for example, their brain unconsciously ramps up their hunger and causes them to eat about 100 more calories. Had our bodies not developed these defense mechanisms, we likely wouldn't have survived the crucible of evolution.

			\hl{This is why fad diets aren't solving the nation's weight problem. It's not information and advice we lack (after all, fad diets work when followed consistently over the long run). It's our inability to persist against the discomfort of hunger---a necessary state for weight loss. Just 3 percent of the people who lose weight in a given year manage to keep it off. Their secret isn't some special food or exercise no one else has. It's their ability to get comfortable with discomfort.} (TCC PDF 137--138)
			
			``I've never believed that people should be doing more or new things. Continuously trying to add more stuff on top of what you're doing and constantly experimenting with shiny new things is almost never the answer. It just adds another layer of stress and complication. I believe people should be doing less and eliminating limiters to progress. It's more effective to modify the behaviors and thought patterns that are keeping you from progressing,'' said Kashey. ``Because your progress is only as good as your most obvious limiter, right?'' (TCC PDF 141--142)
			
		% -----
		\subsubsection{0001 | Easter on junk food}
		% -----
			\label{GM-LIFE-FOOD-0002}
			
			``Processed food is not always junk, but junk is usually processed. I do think junk food is unhealthy, but it's not because sugar is ‘toxic' or any of that nonsense,'' he said. ``It's mainly because it's more calorie dense, less filling, and is more likely to lead someone to overeat and gain weight. And being overweight or obese is one of the largest risk factors for disease.'' (TCC PDF 140)
				
	% ----------------------
	\subsection{Reasoning and Evidence} % *REASONING
	% ----------------------
		\label{GM-LIFE-REASONING}
		
		% -----
		\subsubsection{0001 | Yudkowsky on individual pieces of evidence}
		% -----
			\label{GM-LIFE-REASONING-0001}
			
			Evidence: Until then, he had not felt these extra details as extra burdens. Instead they were corroborative detail, lending verisimilitude to the narrative. Someone presents you with a package of strange ideas, one of which is that universes replicate. Then they present support for the assertion that universes replicate. But this is not support for the package, though it is all told as one story.

			You have to disentangle the details. You have to hold up every one independently, and ask, ``How do we know this detail? (6)
			
		% -----
		\subsubsection{0002 | Yudkowsky on the planning fallacy}
		% -----
			\label{GM-LIFE-REASONING-0002}
			
			Even when asked to make a highly conservative forecast, a prediction that they felt virtually certain that they would fulfill, students' confidence in their time estimates far exceeded their accomplishments.''

			More generally, this phenomenon is known as the ``planning fallacy.'' The planning fallacy is that people think they can plan, ha ha.

			A clue to the underlying problem with the planning algorithm was uncovered by Newby-Clark et al., who found that

			---Asking subjects for their predictions based on realistic ``best guess'' scenarios; and

			---Asking subjects for their hoped-for ``best case'' scenarios...

			...produced indistinguishable results.

			When people are asked for a ``realistic'' scenario, they envision everything going exactly as planned, with no unexpected delays or unforeseen catastrophes---the same vision as their ``best case.''

			Reality, it turns out, usually delivers results somewhat worse than the ``worst case.''

			So there is a fairly reliable way to fix the planning fallacy, if you're doing something broadly similar to a reference class of previous projects. Just ask how long similar projects have taken in the past, without considering any of the special properties of this project. Better yet, ask an experienced outsider how long similar projects have taken. (7)
			
	% ----------------------
	\subsection{Writing} % *WRITING
	% ----------------------
		\label{GM-LIFE-WRITING}
		
		% -----
		\subsubsection{0001 | Yudkowsky on writing for a popular audience}
		% -----
			\label{GM-LIFE-WRITING-0001}
			
			Closely related is the illusion of transparency: We always know what we mean by our words, and so we expect others to know it too. Reading our own writing, the intended interpretation falls easily into place, guided by our knowledge of what we really meant. It's hard to empathize with someone who must interpret blindly, guided only by the words.

			On writing for a popular audience: Combined with the illusion of transparency and self-anchoring, I think this explains a lot about the legendary difficulty most scientists have in communicating with a lay audience---or even communicating with scientists from other disciplines. When I observe failures of explanation, I usually see the explainer taking one step back, when they need to take two or more steps back. Or listeners assume that things should be visible in one step, when they take two or more steps to explain. Both sides act as if they expect very short inferential distances from universal knowledge to any new knowledge. A clear argument has to lay out an inferential pathway, starting from what the audience already knows or accepts. If you don't recurse far enough, you're just talking to yourself.

			If at any point you make a statement without obvious justification in arguments you've previously supported, the audience just thinks you're crazy. (8, 9)
			
	% ----------------------
	\subsection{Spirituality} % *SPIRITUALITY
	% ----------------------
		\label{GM-LIFE-SPIRITUALITY}
		
		% -----
		\subsubsection{0001 | Independent study talk from 2009} % *INDEPENDENTSTUDY
		% -----
			\label{GM-LIFE-SPIRITUALITY-0001}
			
			\textit{[I first read this talk in December of 2009, when I was working as a grader for Independent Study. It was submitted as an assignment in a philosophy class by a woman whose name I cannot remember. It impacted me powerfully as soon as I read it, and it has continued to do that since then.]}

			One of the reasons it was so powerful is that, at the time, I was trying to figure out whether I should marry Alicia. I was looking for spiritual confirmation of the action that I wanted to take, which was to get married to her. After reading this talk, I realized that I had been approaching the problem incorrectly. I hadn't been trying to live very spiritually, very righteously, yet was still seeking guidance from the Lord. This talk prompted me to change several things about what I was going, and led to me receiving a confirmation, one week later, that I should marry her.)

			Our four-year-old son loves to tease by quoting his mother, ``Oh my gosh...Get it off me!...Get it OFF ME!...GET IT OFF ME!'' He finds these statements delightfully humorous because his mother was not awake at the time. You see, I was dreaming I had a huge spider on me. When I am awake, I am not afraid of spiders; yet I was absolutely terrified at that delirious moment. I did not awaken until my own voice woke me up after a series of hollers.

			When we are mentally and physically asleep, we often don't realize the true length of time it takes for experiences to transpire, nor do we always see a connection between our actions in our dreams. We are able to tell when we are physically asleep or awake because the clarity is absent. Similarly, when we are spiritually asleep, our minds are more muddled, for we often don't recognize or appreciate how lasting the positive and negative effects of our thoughts and actions have on ourselves and others, and we sometimes don't see a connection between our individual experiences. My interpretation of one belief of Descartes is that when we have a belief we consider innate, that belief had to be placed, and the only source that could have placed it would be God. So, then, defining what is able to be considered innate becomes my challenge, as we can easily be asleep to truths by considering our own perceptions and philosophies to be innate and from God. How do I know if I am spiritually awake or drowsily walking and talking my way through mortality when following particular cultural, religious, or political views?

			When I was a teenager about twenty years ago, I had a difficult time waking up to my alarm---I unfortunately trained my brain to not register the blaring noise. My poor Dad would come downstairs every morning to inform me that my alarm was going off. I similarly sleep through the reality of Christ's atoning sacrifice just as His three apostles did long ago when I callously accept my weaknesses and sins and don't make more of an effort to change by cultivating His spirit in my heart and home. I sleep through His power to heal just as the guards of the tomb did during His resurrection when I pray simply to check another action off my to-do-list.

			The number of times that the prophets in ancient and modern times warned us to wake up is worth noting\footnote{2 Nephi 1:13--14 and 23; 2 Nephi 4:28; Jacob 3:11; Mosiah 2:38 and 4:5; Alma 4:3, 7:22; Isaiah 51:9; 1 Corinthians 15:34; Romans 13:11.}. When Christ visited the Nephites after his Resurrection, He recognized that they weren't awake to what He was telling them and that they wanted Him to stay with them a little longer. After He asked for the people to kneel in prayer, He groaned within himself and the first thing he said to our Father was, ``Father, I am troubled because of the wickedness of the people...''\footnote{3 Nephi 17:14.}

			Too often I approach Heavenly Father in prayer as if I'm being sent to the principal's office, being too caught in myself, and not trusting in Him enough to know that, at times...he wants me to feel His arms around me. Only when my heart is STILL have I felt our Savior near me--leading my slightly awakened soul to our Father. On that day that Christ suffered for all mankind, He said to his disciples, ``Let not your heart be troubled, neither let it be afraid.''\footnote{John 14:27.} Regarding this, Elder Holland, taught:

			\begin{quotation}
			This may be one of the Savior's commandments that is, in the hearts of otherwise faithful Latter-day Saints, almost universally disobeyed...As concerned as I would be if...one of my children were seriously troubled or disobedient, nevertheless I would be infinitely more devastated if I felt that at such a time that child could not trust me to help, or should feel his or her interest were unimportant to me...I am convinced that none of us can appreciate how deeply it wounds the loving heart of the Savior of the world when he finds that his people do not feel secure in His hands or trust in his commandments...God [and Christ] CANNOT DO OTHER THAN care for us and help us if we will but come [to] them, approaching their throne of grace in meekness and lowliness of heart. It is their nature.\footnote{Holland, Jeffrey. ``Come Unto Me.'' Devotional Speech. Brigham Young University Marriot Center. Provo, Utah. January 17, 1989.}
			\end{quotation}

			Going back to when Christ was praying for the Nephites, verse 3 Nephi 17:15--17 reads:

			\begin{quote}
			Behold, he prayed unto the Father, and the things which he prayed cannot be written...The eye hath never seen, neither hath the ear heard...And no tongue can speak...neither can the hearts of men conceive so great and marvelous things as we both saw and heard Jesus speak; and no one can conceive of the joy which filled our souls at the time we heard him pray for us unto the Father.
			\end{quote}

			Then in verse 20 Jesus says to them, ``Blessed are ye because of your faith. And now behold, my joy is full.'' The joy that was then placed in Jesus was not because of how good the people were living but because, in that moment, they were awakened to the reality of their need for Him, and they had faith in Him. Jesus was troubled because of the wickedness, but in His oneness with God, He ws given a fullness of joy.

			H. Burke Peterson tells of his experience participating, for a second time, in a conference for the inmates at a prison with maximum security located in Southern Arizona.

			\begin{quotation}
			After our proper clearance, we went into a good-sized recreation room where we were to meet with fifty or sixty men. As they came in, I shook hands with many of them. One--I had been his bishop many years earlier. Another--I had been his elders quorum president. Another young man--just a few years earlier I had visited with him as a missionary as we toured one of the missions of Great Britain. These were men who had lost their priesthood, had lost their membership in the Church. [Many had wives, and children who now no longer have a father in the home.] These were men who were there to pay their debt to society.

			Several of them were invited to speak. EACH ONE SPOKE ABOUT THE FEELINGS THEY HAD FOR THEIR SAVIOR, and on repentance, forgiveness, and coming unto Christ. As they spoke, I felt an unusual spirit. The content of their talks was unusual, but beyond that, I was feeling stirrings inside that I hadn't experienced very often: peaceful feelings. As I went home, trying to understand, I turned to 3 Nephi and read the record there of the Beatitudes. I had not understood what I understand now until I had that experience. ``And blessed are all they who do hunger and thirst after righteousness, for they shall be filled with the Holy Ghost.''

			The scripture does not say, blessed are the righteous for they shall be filled with the Holy Ghost, because that is obvious. The scripture says blessed are they who want to be--blessed are they who want MORE THAN ANYTHING ELSE to be [faithful]. I'm beginning to understand, Brothers and Sisters, what that means. It means when I have done the best I can, and when I have not played games with the counsel, the difference between what I can do and what must be done is accomplished because of the grace of Christ.\footnote{Peterson, H. Burke. ``Come Unto Christ Through Your Trials.'' Devotional Speech. Brigham Young University Marriot Center. Provo, Utah. February 16, 1996.}
			\end{quotation}

			I think of this scripture as saying, ``Blessed are those who are awake to their need, who feel their lack, AND who come unto Him.'' When do we realize when we are physically hungry and thirsty? When we are AWAKE...Similarly, we are awakening spiritually when we are hungering and thirsting and seeking to align ourselves with our Father's will. I know that, for me, one of the most valuable forms of being awake to personal revelation comes BEFORE I will be kneeling to pray. It often comes when I am doing such things as taking off my makeup and getting ready for bed, or when I am by myself and thinking about how I have disappointed my Father in Heaven by being unkind to my spouse or children or by thinking too much about worldly cares.

			I have been given a very blessed childhood and life, so the concept of ``having a bad day'' has never felt justifiable to me, personally. There was a moment, though, a few years ago, that was particularly intense (relatively speaking) when I felt a depth of pain that I wasn't aware of. My soul was calling out and yet, simultaneously, His spirit and the love He has for his children was penetrating me so deeply that I remember actually thinking that I didn't want the pain to end, for I didn't want that glimpse of the compassion He has for His children to leave me. When I experienced those feelings simultaneously, I was awakened. As time went on, I began reacting differently to situations that, in the past, I would have critiqued or would have been less apt to giver another the benefit of the doubt. I began listening less to the opinions of those around me (however noble or principle-oriented those opinions may have been) and more to the promptings of the Spirit, and my capacity to distinguish between the two was strengthened. Even though my challenges still remained, He graciously replaced some of their weight with an increased desire to succor those in need. Personal challenges help us to wake up and see things as they really are-and turn away form worldly attractions by valuing our Father's opinion more than man's opinion. I know that, as Jacob testified, the Spirit speaketh of things as they really are and of things as they really will be.\footnote{Jacob 4:13.} Certainly at the point of our sanctification we will want to have had comparable tests to many of our great examples who have gone before us, and we know from where their strength to persist comes.

			I believe we CAN know whether we are spiritually awake by praying and reflecting on the effects that our thoughts and actions have on ourselves and those around us, and then discerning (humbly--by seeking and trusting His spirit rather than our own understanding or biases) what would allow the spirit of our Father and Heaven to remain with us, and what would cause the His spirit to withdraw from us. We can decipher times we are more spiritually awake (than other times) when we are receiving, on a regular basis, promptings and teaching moments that connect the dots between our thoughts, actions, and the effects that we can have on others. I know I am spiritually asleep when my decisions reflect ingratitude and pride, as if I am exempt from the law of justice, or when I feel as though it isn't fair for the sinner (myself) to be miserable rather than being humble and faithful in changing and exercising the power of the atonement and finding that joy once again. I don't know of anything that deprives us of this strength and joy more than being asleep to the commandments of our Father. Waking up to the teachings of Christ through repentance is one of the wonderful realities of His plan for us.
			
		% -----
		\subsubsection{0002 | Bobpa's testimony} % *BOBPA
		% -----
			\label{GM-LIFE-SPIRITUALITY-0002}
			
			\textit{[In late 2012, Tina sent her father's written testimony to all of her children and their spouses. His name is Robert Fletcher, but all of the grandkids call him Bobpa, and I call him that as well.]}
			
			I am writing this not because I find fault with the way beliefs are stated by many members of the LDS church. It is primarily written for young people for whom the Internet is a way of life from which all points of view are freely expressed. As never before, information about the church is thrust in front of them for the purpose of demonstrating its falsehood. I would like to describe my own way of dealing with troublesome historical ``facts''. I will start by listing some of my key beliefs. They are not significantly different from those of most LDS members, even though my way of describing them is sometimes unconventional:

			I believe that each of us has a spirit as well as a body as manifested by our awareness and our ability to make decisions that are carried out by our body (free agency). (Scientists are striving to show that this awareness comes physically through the operation of the brain, so far unsuccessfully. Nevertheless, all the laws of a country are based on the assumption that we have free will.)

			Thus, I believe in the reality of spirit.

			It doesn't seem to be too much of a stretch to believe that our spirits existed independent of our bodies before we were born and will exist after our bodies die.

			I believe in God as the Supreme Being who has provided a plan for the possibility of giving us continuous and ultimate ``joy.''

			I believe in Jesus Christ as His Son who has implemented God's plan, and taught us His gospel.

			I believe in the Holy Ghost that influences our spirit.

			I believe that Joseph Smith was a prophet who was directed by God to form His church.

			I believe that true prophets are described in the Bible and the Book of Mormon who were inspired by the Holy Ghost.

			I believe that one of the main purposes for our mortal existence is to develop the ability to live a Christ-like life in the presence of mortal temptations to do otherwise, with no memory of our previous spiritual life and with little physical evidence of God or other spiritual beings to guide us, being left to be guided only by the Holy Ghost.

			I believe that our Spirit as developed by these choices will find an opportunity in the here-after uniquely suited to our mortal development thanks to the eternal plan of Jesus Christ (Atonement).

			Up to this point I think my above stated beliefs do not differ from, but of course do not include all, the beliefs of the LDS church. What follows is my own view of some of these beliefs. 

			I will not try to distinguish between the influences of various forms of divine spirit: the influence of the Holy Ghost, the Spirit of God, the Light of Christ, the Lord's spirit, etc., and I will refer to all as the influence of the Holy Ghost, which is His specific mission. I believe that all spiritual manifestations are aspects of that influence, including impressions, burning of one's bosom, words of guidance received in one's mind, or visions that may be so vivid as to seem physical (but often described as being seen with spiritual eyes). By physical, I mean obeying the scientific laws of nature, such as the traveling of photons through space as emitted by a source such as the Sun or a burning bush, or as reflected off physical objects, or as perceived by the retina of our physical eyes and recorded by our brains. Visions may not be physical, but be spiritual influences impressed directly on our spirit. This does not rule out the existence of beings with spatial and time locality, such as angels and resurrected beings, but they clearly have properties and powers foreign to our mortal knowledge.

			I believe that the influence of the Holy Ghost is most often not precise or all inclusive. It is a subtle influence that combines with the reasoning power of our brains to lead us to a decision that is favored by God compared to other decisions we are considering that are not so favored. When Nephi was impressed to cut off Laban's head, this impression did not mean that God condones decapitation under other circumstances. Our brains are influenced by the memories of all our past experiences, including the teachings of our parents and schools and by the influence of the culture(s) in which we have existed. Probably most of our decisions are a combination of reasoning impacting on that collective experience without much current influence of the Holy Ghost. Our decisions are often strongly influenced by the mortal need to satisfy our mortal appetites. Even the greatest prophets had these influences. These affected the way in which they expressed their prophetic utterances.

			When I read the scriptures, particularly the Old Testament, it's hard to believe that a loving God was as concerned as were the prophets in overcoming enemies. Thus, Saul was commanded through the prophet Samuel to slay all the Amalekites (\hyperref[1SAMUEL-15]{1 Samuel 15}) and was rebuked when he didn't slay the King and some of the sheep. Saul was living at a time when warring and killing was part of their culture. So within that culture God (presumably speaking through Samuel) directed Saul to do something that he would not condone in other times in order to accomplish some greater long-term purpose. We cannot conclude that because Samuel was impressed at that time to direct a mass slaying is that this is what God always wants. This is an illustration that the guidance of the Holy Ghost tends not to be precise nor all inclusive. One of God's greatest concerns in those days was to develop recognition that there was one personal God that was of much greater importance to their mortality than were the idols and pagan gods that were commonly worshiped. Through all of their other concerns, the prophets of the Old Testament managed to get guidance to that end. But their preserved descriptions of their spiritual experiences in the scriptures retain the influences of the culture in which they were embedded.

			As we try to understand the scriptures, we should also keep in mind that descriptions of the scriptural events are influenced both by the passage of time and by the description by the observers in terms that they believed. Thus, the stories in my own family first by me and my wife and then by my children have grown in ways that make them more interesting, probably at the expense of accuracy. I do not doubt that same kind of thing happened to stories in the scriptures.

			I believe Joseph Smith was a prophet in these Latter Days. He existed in a culture that believed in transcendence of spiritual over physical laws (miracles). There were many in his surroundings who claimed they had been visited personally by Jesus Christ. Many believed in buried treasure, peep stones, and divining rods. Thus, it did not surprise many when Joseph Smith claimed to have visions of buried gold plates and was impressed that he should find them and translate them. He was led to do so, but there is no evidence that he even had to look at the plates to translate them. With little formal education, and by remarkable spiritual impressions, he was able to dictate the Book of Mormon in a period of no more than three months. Eleven witnesses had miraculous visions of the existence of the plates, which they never denied. It is almost impossible to believe that Joseph with his lack of formal education could have created the translation without spiritual guidance. When the book was printed and distributed, people believed it was physical evidence for what he claimed. It became (and still is) the greatest influence on its readers to convince them that Joseph Smith was a prophet and that the church he established was guided by God.

			However, I don't necessarily believe that everything in the book is historically accurate. He created sentences in his mind that impressed him as good, and that was what he dictated (see \hyperref[DC9]{DC 9}). The result is a book with remarkable spiritual insight that can lead a person to live a Christ-like life. But I will not be bothered if there is physical evidence that the history described therein is unlikely to be exactly as described. Joseph's mother recalled that Joseph regaled the family with stories of ancient America even before he received the plates. I would guess that these ``pre-visions'' influenced him as he translated.

			The same process presumably applied to the translation of the Book of Abraham in the Pearl of Great Price. Here the papyri from which he derived his impressions were clearly physical. They still exist in the possession of the Church. Copies have been made available to scholars outside the church. These scholars have mostly believed that the papyri represent Egyptian hieroglyphics from the Egyptian Book of the Dead. Although there are LDS scholars (eg, Hugh Nibley) that dispute this, it is not hard for me to believe that Joseph Smith received impressions from the papyri, not literal translations, which led him to the creation of the Book of Abraham, a book with great insight into spiritual matters. Joseph used his own definition of truth: ``Whatsoever thing is good is...true.'' (\hyperref[MORONI-10-6]{Moroni 10:6})

			The danger of deciding that some part of the scriptures is the influence of culture and not God is to allow ourselves exceptions to some of God's commandments if they do not suit us. But I believe this is a part of God's plan. We are to decide for ourselves what is right and good by heeding to the Holy Ghost. This is what gives us a testimony of Joseph Smith and the church that God directed him to establish. But we need to recognize that rationalization is also a temptation that we need to resist if it leads us to non-Christ-like decisions. One of the purposes of continuous revelation is to restate truths and practices as the culture changes. There is a greater need for this today as never before with the rapid changes in our culture brought about by technology.

			I believe that the normal description of spirit and spiritual beings are simple models that serve most people very well, guiding them to Christ-like living. We visualize God and angels as physical beings not unlike ourselves. But clearly they are endowed with properties that we can't understand or visualize. When I pray to God, I don't visualize Him in any particular location, but as an exalted being, listening with loving understanding, even when my prayer is unvoiced. I don't try to understand how He does this simultaneously with the millions of other people that might be praying at the same time. The simple model is appropriate to make my prayers meaningful to me and I hope to Him.

			I tend to think of the Holy Ghost (or one of his spiritual helpers) as a person transmitting a message to me by something like mental telepathy. This is the model that I think most LDS members have. The reality is probably quite different. I am unable to even speculate how the spiritual influence really works. But we do not need a better model in order to have our decisions be influenced by Him. The same applies to visions. If a person feels comfortable with the belief that the visit of a heavenly person is a physical reality with physical photons emitted, so be it. Whatever model works to convince one of the reality of the visitation by the testimony of the spirit is good. In fact, I believe that the attempt to attach physical laws of nature to spiritual experience is often counter-productive; the effort to do so can interfere with the benefits of the spiritual experience, e.g., in personal prayers or studying the scriptures. Some scholars call this process transcendence (belief of spiritual truths over logic or physical laws of nature). Whatever the model or transcendence, if the witness of the Holy Ghost generates faith in Jesus Christ and his gospel and motivates a person to live that gospel, I believe that helps him/her accomplish the purpose for which God gave him/her this mortal experience.

			I believe that the Holy Ghost is available to all people if they look for it and adjust their lives by its guidance. Some LDS members describe this as the Light of Christ. I think that God is concerned for all his children on earth and is pleased whenever they are influenced to live a Christ-like life, no matter what religion they adhere to. I believe that God loves all his children on this earth and always makes available His spirit to guide them. But it is up to each individual to live a life in which the spirit can be heard. I pray nightly for some of my own grandchildren who have become inactive in the church that they will listen to His spirit and be guided into a good life.

			In conclusion, my advice to a youth is to decide on the way of life that seems good to you, recognizing that you have a spirit that probably has an eternal existence. Do not get hung up on perceived deficiencies of prophets or other leaders, but pay attention to the main thrust of their messages. My own conclusion is that following the teachings and practices of the Church of Jesus Christ of Latter Day Saints is the best way I know that helps me approach a Christ-like life, although for me I have a long way to go. One of the greatest external evidences of this is to see the change in people's lives when they witness the Spirit of the Holy Ghost and join with fellow members of the church in striving to become more Christ-like. I have received guidance of the spirit in my life that has helped me make choices and given me assurance and peace that His gospel is true, and that He is pleased when I try to live it.
			
		% -----
		\subsubsection{0003 | Plato's allegory of the cave} % *CAVE
		% -----
			\label{GM-LIFE-SPIRITUALITY-0003}
			
			\textit{[I have copied and pasted this text from} Republic book VII \textit{514a--520a, from the Hackett complete works edition.]}
		
			Next, I said, compare the effect of education and of the lack of it on our nature to an experience like this: Imagine human beings living in an underground, cavelike dwelling, with an entrance a long way up, which is both open to the light and as wide as the cave itself. They've been there since childhood, fixed in the same place, with their necks and legs fettered, able to see only in front of them, because their bonds prevent them from turning their heads around. Light is provided by a fire burning far above and behind them. Also behind them, but on higher ground, there is a path stretching between them and the fire. Imagine that along this path a low wall has been built, like the screen in front of puppeteers above which they show their puppets. 

			I'm imagining it. 

			Then also imagine that there are people along the wall, carrying all kinds of artifacts that project above it---statues of people and other animals, made out of stone, wood, and every material. And, as you'd expect, some of the carriers are talking, and some are silent. 

			It's a strange image you're describing, and strange prisoners. 

			They're like us. Do you suppose, first of all, that these prisoners see anything of themselves and one another besides the shadows that the fire casts on the wall in front of them? 

			How could they, if they have to keep their heads motionless throughout life? 

			What about the things being carried along the wall? Isn't the same true of them? 

			Of course. 

			And if they could talk to one another, don't you think they'd suppose that the names they used applied to the things they see passing before them? 

			They'd have to. 

			And what if their prison also had an echo from the wall facing them? Don't you think they'd believe that the shadows passing in front of them were talking whenever one of the carriers passing along the wall was doing so? 

			I certainly do. 

			Then the prisoners would in every way believe that the truth is nothing other than the shadows of those artifacts. 

			They must surely believe that. 

			Consider, then, what being released from their bonds and cured of their ignorance would naturally be like, if something like this came to pass.2 When one of them was freed and suddenly compelled to stand up, turn his head, walk, and look up toward the light, he'd be pained and dazzled and unable to see the things whose shadows he'd seen before. What do you think he'd say, if we told him that what he'd seen before was inconsequential, but that now---because he is a bit closer to the things that are and is turned towards things that are more---he sees more correctly? Or, to put it another way, if we pointed to each of the things passing by, asked him what each of them is, and compelled him to answer, don't you think he'd be at a loss and that he'd believe that the things he saw earlier were truer than the ones he was now being shown? 

			Much truer. 

			And if someone compelled him to look at the light itself, wouldn't his eyes hurt, and wouldn't he turn around and flee towards the things he's able to see, believing that they're really clearer than the ones he's being shown? 

			He would. 

			And if someone dragged him away from there by force, up the rough, steep path, and didn't let him go until he had dragged him into the sunlight, wouldn't he be pained and irritated at being treated that way? And when he came into the light, with the sun filling his eyes, wouldn't he be unable to see a single one of the things now said to be true? 

			He would be unable to see them, at least at first. 

			I suppose, then, that he'd need time to get adjusted before he could see things in the world above. At first, he'd see shadows most easily, then images of men and other things in water, then the things themselves. Of these, he'd be able to study the things in the sky and the sky itself more easily at night, looking at the light of the stars and the moon, than during the day, looking at the sun and the light of the sun. 

			Of course. 

			Finally, I suppose, he'd be able to see the sun, not images of it in water or some alien place, but the sun itself, in its own place, and be able to study it. 

			Necessarily so. 

			And at this point he would infer and conclude that the sun provides the seasons and the years, governs everything in the visible world, and is in some way the cause of all the things that he used to see. 

			It's clear that would be his next step. What about when he reminds himself of his first dwelling place, his fellow prisoners, and what passed for wisdom there? Don't you think that he'd count himself happy for the change and pity the others? 

			Certainly. 

			And if there had been any honors, praises, or prizes among them for the one who was sharpest at identifying the shadows as they passed by and who best remembered which usually came earlier, which later, and which simultaneously, and who could thus best divine the future, do you think that our man would desire these rewards or envy those among the prisoners who were honored and held power? Instead, wouldn't he feel, with Homer, that he'd much prefer to ``work the earth as a serf to another, one without possessions,'' and go through any sufferings, rather than share their opinions and live as they do? 

			I suppose he would rather suffer anything than live like that. 

			Consider this too. If this man went down into the cave again and sat down in his same seat, wouldn't his eyes---coming suddenly out of the sun like that---be filled with darkness? 

			They certainly would. 

			And before his eyes had recovered---and the adjustment would not be quick---while his vision was still dim, if he had to compete again with the perpetual prisoners in recognizing the shadows, wouldn't he invite ridicule? Wouldn't it be said of him that he'd returned from his upward journey with his eyesight ruined and that it isn't worthwhile even to try to travel upward? And, as for anyone who tried to free them and lead them upward, if they could somehow get their hands on him, wouldn't they kill him? 

			They certainly would. 

			This whole image, Glaucon, must be fitted together with what we said before. The visible realm should be likened to the prison dwelling, and the light of the fire inside it to the power of the sun. And if you interpret the upward journey and the study of things above as the upward journey of the soul to the intelligible realm, you'll grasp what I hope to convey, since that is what you wanted to hear about. Whether it's true or not, only the god knows. But this is how I see it: In the knowable realm, the form of the good is the last thing to be seen, and it is reached only with difficulty. Once one has seen it, however, one must conclude that it is the cause of all that is correct and beautiful in anything, that it produces both light and its source in the visible realm, and that in the intelligible realm it controls and provides truth and understanding, so that anyone who is to act sensibly in private or public must see it. 

			I have the same thought, at least as far as I'm able. 

			Come, then, share with me this thought also: It isn't surprising that the ones who get to this point are unwilling to occupy themselves with human affairs and that their souls are always pressing upwards, eager to spend their time above, for, after all, this is surely what we'd expect, if indeed things fit the image I described before. 

			It is. 

			What about what happens when someone turns from divine study to the evils of human life? Do you think it's surprising, since his sight is still dim, and he hasn't yet become accustomed to the darkness around him, that he behaves awkwardly and appears completely ridiculous if he's compelled, either in the courts or elsewhere, to contend about the shadows of justice or the statues of which they are the shadows and to dispute about the way these things are understood by people who have never seen justice itself? 

			That's not surprising at all. 

			No, it isn't. But anyone with any understanding would remember that the eyes may be confused in two ways and from two causes, namely, when they've come from the light into the darkness and when they've come from the darkness into the light. Realizing that the same applies to the soul, when someone sees a soul disturbed and unable to see something, he won't laugh mindlessly, but he'll take into consideration whether it has come from a brighter life and is dimmed through not having yet become accustomed to the dark or whether it has come from greater ignorance into greater light and is dazzled by the increased brilliance. Then he'll declare the first soul happy in its experience and life, and he'll pity the latter---but even if he chose to make fun of it, at least he'd be less ridiculous than if he laughed at a soul that has come from the light above. 

			What you say is very reasonable. 

			If that's true, then here's what we must think about these matters: Education isn't what some people declare it to be, namely, putting knowledge into souls that lack it, like putting sight into blind eyes. 

			They do say that. 

			But our present discussion, on the other hand, shows that the power to learn is present in everyone's soul and that the instrument with which each learns is like an eye that cannot be turned around from darkness to light without turning the whole body. This instrument cannot be turned around from that which is coming into being without turning the whole soul until it is able to study that which is and the brightest thing that is, namely, the one we call the good. Isn't that right? 

			Yes. 

			Then education is the craft concerned with doing this very thing, this turning around, and with how the soul can most easily and effectively be made to do it. It isn't the craft of putting sight into the soul. Education takes for granted that sight is there but that it isn't turned the right way or looking where it ought to look, and it tries to redirect it appropriately. 

			So it seems. 

			Now, it looks as though the other so-called virtues of the soul are akin to those of the body, for they really aren't there beforehand but are added later by habit and practice. However, the virtue of reason seems to belong above all to something more divine, which never loses its power but is either useful and beneficial or useless and harmful, depending on the way it is turned. Or have you never noticed this about people who are said to be vicious but clever, how keen the vision of their little souls is and how sharply it distinguishes the things it is turned towards? This shows that its sight isn't inferior but rather is forced to serve evil ends, so that the sharper it sees, the more evil it accomplishes. 

			Absolutely. 

			However, if a nature of this sort had been hammered at from childhood and freed from the bonds of kinship with becoming, which have been fastened to it by feasting, greed, and other such pleasures and which, like leaden weights, pull its vision downwards---if, being rid of these, it turned to look at true things, then I say that the same soul of the same person would see these most sharply, just as it now does the things it is presently turned towards. 

			Probably so. 

			And what about the uneducated who have no experience of truth? Isn't it likely---indeed, doesn't it follow necessarily from what was said before---that they will never adequately govern a city? But neither would those who've been allowed to spend their whole lives being educated. The former would fail because they don't have a single goal at which all their actions, public and private, inevitably aim; the latter would fail because they'd refuse to act, thinking that they had settled while still alive in the faraway Isles of the Blessed. 

			That's true. 

			It is our task as founders, then, to compel the best natures to reach the study we said before is the most important, namely, to make the ascent and see the good. But when they've made it and looked sufficiently, we mustn't allow them to do what they're allowed to do today. 

			What's that? 

			To stay there and refuse to go down again to the prisoners in the cave and share their labors and honors, whether they are of less worth or of greater. 

			Then are we to do them an injustice by making them live a worse life when they could live a better one? 

			You are forgetting again that it isn't the law's concern to make any one class in the city outstandingly happy but to contrive to spread happiness throughout the city by bringing the citizens into harmony with each other through persuasion or compulsion and by making them share with each other the benefits that each class can confer on the community.4 The law produces such people in the city, not in order to allow them to turn in whatever direction they want, but to make use of them to bind the city together.